\section{Capital Architecture}

The Alliance funds itself in three layers, each with a distinct source,
mandate, and time horizon. The layered design protects every individual
member from being starved during a slow period elsewhere in the network,
and protects the Alliance's coordination layer from any single member's
funding cycle.

\subsection{Layer 1 --- Substrate and coordination capital}

\textbf{Source.} Strategic investors with multi-cycle horizons. Family
offices and sovereign-adjacent capital with discretionary mandates.
Founders' personal equity that has been re-invested into the Alliance
holding rather than monetised at prior exits.

\textbf{Mandate.} Fund (i) the proprietary L2 substrate's continued
development, (ii) the patent estate's prosecution and defence, (iii) the
coordination layer's operating costs (governance, legal, finance, treasury,
compliance, secops), and (iv) reserve runway against any single member's
shortfall.

\textbf{Time horizon.} Ten-year minimum lockup. Returns are distributed
through participation in the Alliance's exits rather than current cash
distribution.

\textbf{Custodian.} The W3A Foundation\ parent holding (Delaware)
with subsidiary structures by jurisdiction. All Layer-1 capital sits in
this vehicle; member-level vehicles draw from it on negotiated terms.

\subsection{Layer 2 --- Member operating capital}

\textbf{Source.} Member-level rounds underwritten primarily by the
Alliance treasury, supplemented by external venture or strategic capital
where it adds value (e.g., a strategic enterprise customer who wants
equity exposure to a specific member, or a sovereign development bank
funding a regional expansion).

\textbf{Mandate.} Fund member-level operating costs, marketing, customer
acquisition, regional licensing, and product development specific to the
member. Each member is responsible for its own P\&L; the Alliance does
not subsidise underperformance indefinitely (\S 8.3 governance).

\textbf{Time horizon.} Standard venture cycle (5--7 years to liquidity
event or sustained profitability), but with the optionality of an
Alliance-internal recapitalisation rather than a forced external exit
when external markets are uncooperative.

\subsection{Layer 3 --- Acquisition and spawn capital}

\textbf{Source.} Operating cash flow from Phase-Two financial super-app
revenue, supplemented by Alliance treasury and by structured-debt
facilities collateralised against the Alliance's IP estate and
substrate-revenue stream.

\textbf{Mandate.} Fund Phase-Three transactions (\S 5): acquisition
consideration, spawn-vehicle initial capitalisation, integration cost,
and post-close acceleration capital.

\textbf{Time horizon.} Transaction-specific. Each Phase-Three transaction
carries its own thesis with a defined IRR target and a defined exit
horizon. The Alliance treasury manages the portfolio of these
transactions as a single book.

\subsection{Cross-layer rules}

\textbf{No round-tripping.} Layer-1 capital does not fund Layer-3
acquisitions directly --- it funds substrate, not member operations. Layer-3
capital is sourced from Layer-2 cash flow or external acquisition
facilities, not from substrate reserves.

\textbf{No cross-subsidisation by default.} Member A's operating cash
flow does not fund member B's losses unless the coordination layer
explicitly approves a member-stabilisation distribution under documented
criteria.

\textbf{All capital is patient.} The Alliance does not take capital that
demands liquidity inside 24 months. Members who require shorter-cycle
capital raise it themselves at the member level, against their own P\&L,
without committing the Alliance's larger capital pool.

\textbf{Documented IRR targets per layer.} Layer 1 targets 25--40\%
gross IRR over the ten-year horizon, weighted heavily by terminal
enterprise value. Layer 2 targets 30--60\% per individual member round
(standard venture math) with a portfolio expectation of 25\% gross IRR.
Layer 3 targets 25--35\% per individual transaction with a portfolio
expectation of 30\% gross IRR.

\subsection{Treasury policy}

The Alliance treasury is operated under documented policy that
prioritises capital preservation over yield optimisation. Idle cash is
held in money-market instruments, short-duration treasury bills, and
investment-grade corporate paper. Crypto-asset treasury is denominated
primarily in dollar-pegged stablecoins held in segregated custody under
the W3A MPC + KMS posture. The treasury does not speculate on the price
of the Alliance's own native tokens; member tokens issued under Phase-
Three transactions are vested in the treasury on the same schedule as
employee equity.
